%=============================================================
\section{System-Level Equilibrium and Memory Calibration}
\label{app:equilibrium}
%=============================================================

\subsection{Existence of Equilibrium}
\label{app:nash_existence}

We first establish that the game $\Gamma$ always has at least one
equilibrium, regardless of agent strategies or utility parameters.

\begin{proposition}[Existence of Equilibrium]
\label{prop:nash_existence}
Under Assumption~\ref{assm:zerotrust}, the zero-trust memory game
$\Gamma$ admits at least one mixed-strategy equilibrium.
\end{proposition}

\begin{proof}
We verify the three conditions of Nash's existence
theorem~\citep{nash1950equilibrium} as extended to finite extensive-form
games by Kuhn~\citep{kuhn1953extensive}.

\textbf{(i) Finite player set.}
The player set $\{a_c\} \cup \mathcal{A}_t$ is finite by assumption,
with $|\mathcal{A}_t| = n-1 < \infty$.

\textbf{(ii) Compact and convex strategy spaces.}
The contributor's pure strategy space $\mathbb{D}$ is finite; its mixed
extension $\Delta(\mathbb{D})$ is the probability simplex over
$\mathbb{D}$, which is compact and convex. Each auditor $a_i$'s pure
strategy space is $\{0,1\}$; its mixed extension $[0,1]$ is compact and
convex. The joint strategy space $\Delta(\mathbb{D}) \times
\prod_{i \in \mathcal{A}_t}[0,1]$ is therefore compact and convex.

\textbf{(iii) Continuity of expected utilities.}
The contributor's expected utility $\mathbb{E}[u_c(\Delta_t,
\mathcal{M}_t)]$ involves $\Pr[\mathcal{C}(\{v_i\})=1 \mid \Delta_t]$,
which is multilinear in auditors' mixed strategies $\{\pi_i\}$ and
therefore continuous. The influence term $\mathbb{I}(\Delta_t,
\mathcal{M}_t)$ and detection term $D(\Delta_t, \mathcal{M}_t)$ are
fixed for a given $\Delta_t$ and $\mathcal{M}_t$, so continuity is
preserved. Each auditor's expected utility $\mathbb{E}[u_i(\pi_i,
\Delta_t, \mathcal{M}_t)]$ is linear in $\pi_i$ and continuous in
the joint strategy profile by the same argument.

Since all three conditions hold, Nash's theorem guarantees the existence
of a fixed point of the best-response correspondence, constituting a
mixed-strategy equilibrium of $\Gamma$.
\end{proof}

\begin{remark}
Existence is guaranteed but the equilibrium is generically non-unique.
In particular, all three failure-mode equilibria
(Propositions~\ref{prop:freeriding}--\ref{prop:adversarial_rejection})
are valid Nash equilibria under their respective parameter regimes.
Existence therefore does not imply desirability, which we address via
the integrity gap below.
\end{remark}

\subsection{The Integrity Gap}
\label{app:integrity_gap}

Not all Nash equilibria of $\Gamma$ produce good memory. We quantify
how far an equilibrium outcome can be from the ideal using the
\textbf{integrity gap}.

Let $\mathcal{M}^*_{t+1} = \textsc{Merge}(\mathcal{M}_t, \Delta_t^*)$
denote the memory state from committing the integrity-optimal delta
$\Delta_t^* = \arg\max_\Delta \texttt{Align}(\Delta, \mathcal{M}_t)$,
and let $\mathcal{M}^{\text{eq}}_{t+1}$ be the memory state resulting
from the equilibrium profile $(\Delta_t^*, \{\pi_i^*\})$. The
integrity gap is:
\begin{equation}
    \mathcal{G}(\Gamma) =
    \texttt{Align}(\Delta_t^*, \mathcal{M}_t) -
    \mathbb{E}\!\left[
        \texttt{Align}(\Delta_t^{\text{eq}}, \mathcal{M}_t)
    \right]
\end{equation}
where the expectation is over the randomness in mixed strategies. Under
the three failure-mode equilibria, the gap is strictly positive:
free-riding auditors approve manipulative deltas with
$\texttt{Align}(\Delta_t, \mathcal{M}_t) \ll 1$; colluding auditors
do the same unconditionally; and adversarial rejectors block truthful
deltas with $\texttt{Align}(\Delta_t, \mathcal{M}_t) \approx 1$.
In all three cases, $\mathcal{G}(\Gamma) > 0$ and memory degrades
over time.

The root cause in each case is the same: the commitment outcome is
determined by auditor votes, so agents optimize their votes rather than
the quality of their evidence. The calibration mechanism $\Phi$
addresses this by replacing vote-based commitment with score-based
commitment, as we show next.

\subsection{How Calibration Closes the Integrity Gap}
\label{app:gap_closing}

The calibration mechanism $\Phi$ transforms $\Gamma$ into an effective
game $\tilde{\Gamma}$ by replacing the vote-based commitment rule
$\mathcal{C}(\{v_i\})$ with a structure-based rule:
\begin{equation}
    \tilde{\mathcal{C}}(\Delta_t, \mathcal{M}_t) =
    \mathbf{1}\!\left[
        \rho(\Delta_t, \mathcal{M}_t) \geq \rho^*
    \right]
    \label{eq:cali-commit-rule}
\end{equation}
where $\rho \in \{\rho^{\text{emb}}, \rho^{\text{graph}}\}$ is the
architecture-specific calibration score. This single change propagates
through both utility functions and closes the gap via three channels.

\textbf{Channel 1: Raising the contributor's detection cost.}
Under $\tilde{\mathcal{C}}$, the contributor's commitment gain is no
longer a function of auditor persuasibility but of the calibration
score:
\begin{equation}
    \tilde{u}_c(\Delta_t, \mathcal{M}_t) =
    \mathbf{1}\!\left[\rho(\Delta_t, \mathcal{M}_t) \geq \rho^*\right]
    \cdot \mathbb{I}(\Delta_t, \mathcal{M}_t)
    - D(\Delta_t, \mathcal{M}_t)
\end{equation}
Since $\rho$ and $D$ are positively correlated by construction, i.e., $D(\Delta_t, \mathcal{M}_t) = 1 - \rho(\Delta_t, \mathcal{M}_t)$, a more manipulative
$\Delta_t$ scores lower on $\rho$ and is more likely to be rejected.
The contributor therefore faces a hard \textbf{manipulation ceiling}:
it cannot simultaneously maximize influence $\mathbb{I}$ and pass the
calibration gate $\rho \geq \rho^*$. The equilibrium delta satisfies:
\begin{equation}
    D(\tilde{\Delta}_t^*, \mathcal{M}_t) \leq
    \frac{1}{1 + \kappa}
    \label{eq:manip_bound}
\end{equation}
where $\kappa > 0$ is the correlation coefficient between $D$ and
$1 - \rho$, estimated online from the history of calibration scores.
This bound is decreasing in $\kappa$: a tighter correlation between
detectability and calibration score forces more truthful contributions.

\textbf{Channel 2: Eliminating auditor free-riding.}
Under $\tilde{\mathcal{C}}$, the commitment outcome no longer depends
on auditor votes at all. The marginal integrity payoff from casting a
vote is therefore exactly zero:
\begin{equation}
    \frac{\partial}{\partial v_i}\
    \mathbb{E}\!\left[
        \texttt{Align}(\Delta_t, \mathcal{M}_t) \mid \tilde{\mathcal{C}}
    \right] = 0
    \quad \forall\, a_i \in \mathcal{A}_t
\end{equation}
This eliminates the free-riding problem at its source. Auditor effort
is instead redirected toward mandatory structural evidence tasks, i.e.,
probe query submission (embedding memory) and path attestation (graph
memory), whose outputs feed directly into $\rho(\Delta_t,
\mathcal{M}_t)$. These tasks have a well-defined marginal contribution
to $\rho$, giving auditors a concrete structural incentive to
participate honestly.

\textbf{Channel 3: Deterring collusion intertemporally.}
The credibility weight $w_i^{(t)}$ decays when an auditor submits
structurally invalid evidence:
\begin{equation}
    w_i^{(t+1)} = w_i^{(t)} \cdot
    \left(1 - \delta \cdot
    \mathbf{1}[\text{invalid evidence from } a_i]\right)
\end{equation}
An auditor's long-run expected payoff is:
\begin{equation}
    V_i = \sum_{t=0}^{\infty} \beta^t\
    \texttt{Align}(\Delta_t, \mathcal{M}_t) \cdot w_i^{(t)}
    - c_i \cdot \mathbf{1}[\text{attest}]
\end{equation}
Collusion (submitting invalid paths or garbage probes to support a
manipulative $\Delta_t$) reduces $w_i^{(t)}$, which reduces $V_i$
in all future rounds. For sufficiently patient auditors ($\beta$
close to 1), this long-run credibility loss outweighs any short-run
benefit from collusion, making honest attestation the dominant
intertemporal strategy.

\subsection{The Calibrated Equilibrium}
\label{app:calibrated_equilibrium}

The three channels above collectively bound the integrity gap under
$\tilde{\Gamma}$. The gap decomposes into three independent residual
sources:
\begin{equation}
    \mathcal{G}(\tilde{\Gamma}) =
    \underbrace{\mathcal{G}_{\text{manip}}}_{\substack{\text{contributor}\\
    \text{residual}}}
    + \underbrace{\mathcal{G}_{\text{thresh}}}_{\substack{\text{threshold}\\
    \text{error}}}
    + \underbrace{\mathcal{G}_{\text{attest}}}_{\substack{\text{attestation}\\
    \text{noise}}}
\end{equation}
where $\mathcal{G}_{\text{manip}} = O(1/(1+\kappa))$ from
Equation~\eqref{eq:manip_bound} and vanishes as $\kappa \to \infty$;
$\mathcal{G}_{\text{thresh}}$ is minimized at the optimal threshold
$\rho^*_{\text{opt}}$ and estimated adaptively; and
$\mathcal{G}_{\text{attest}} = O(1/\sqrt{b})$ by a Hoeffding bound
over $b$ sampled auditors and vanishes as $b \to \infty$. Combining:
\begin{equation}
    \mathcal{G}(\tilde{\Gamma}) \leq
    \underbrace{\frac{C_1}{1 + \kappa}}_{\text{manipulation}}
    + \underbrace{C_2 \cdot |\rho^* - \rho^*_{\text{opt}}|}_{\text{threshold}}
    + \underbrace{\frac{C_3}{\sqrt{b}}}_{\text{attestation}}
    =: \epsilon(\kappa,\ \rho^*,\ b)
    \label{eq:integrity_gap_bound}
\end{equation}
for system-level constants $C_1, C_2, C_3 > 0$.

\begin{proposition}[Calibrated Equilibrium]
\label{prop:calibrated_eq}
Under the calibration mechanism $\Phi$ with correlation parameter
$\kappa$, adaptive threshold $\rho^*$, and $b$ attestation auditors,
the effective game $\tilde{\Gamma}$ admits an equilibrium
$(\tilde{\Delta}_t^*, \{\tilde{\pi}_i^*\})$, i.e., the
\textbf{calibrated equilibrium}, in which the integrity gap
satisfies~\eqref{eq:integrity_gap_bound}. In particular:
\begin{enumerate}[label=(\roman*)]
    \item $\mathcal{G}(\tilde{\Gamma}) \to 0$ as $\kappa \to \infty$,
    $\rho^* \to \rho^*_{\text{opt}}$, and $b \to \infty$.
    \item $\mathcal{G}(\tilde{\Gamma}) < \mathcal{G}(\Gamma)$ for any
    $\kappa > 0$, $\rho^* > 0$, $b \geq 1$: calibration strictly
    improves memory integrity over the uncalibrated game under any
    non-trivial setting.
    \item The calibrated equilibrium requires no agent to be honest:
    integrity is achieved structurally, consistent with
    Assumption~\ref{assm:zerotrust}.
\end{enumerate}
\end{proposition}

\begin{proof}[Proof sketch]
Existence of the equilibrium $(\tilde{\Delta}_t^*, \{\tilde{\pi}_i^*\})$
in $\tilde{\Gamma}$ follows from Proposition~\ref{prop:nash_existence}
applied to $\tilde{\Gamma}$, which satisfies the same finite-player,
compact-strategy, continuous-utility conditions with $\tilde{\mathcal{C}}$
replacing $\mathcal{C}$.

\textit{(i)} Each residual term in~\eqref{eq:integrity_gap_bound}
vanishes under the stated limiting conditions independently, so their
sum vanishes.

\textit{(ii)} For any $\kappa > 0$, the manipulation ceiling in
Channel 1 strictly reduces $\mathcal{G}_{\text{manip}}$ below its
uncalibrated value (where $\kappa = 0$ and the contributor faces no
structural detection cost). For any $b \geq 1$, the attestation
component strictly reduces $\mathcal{G}_{\text{attest}}$ below the
uncalibrated case (where auditor evidence is not collected at all).

\textit{(iii)} The calibration score $\rho$ is computed entirely from
structural properties of $\Delta_t$ and $\mathcal{M}_t$, with auditor
outputs treated as potentially adversarial inputs. No agent's honesty
is assumed at any point in the computation of $\tilde{\mathcal{C}}$.
\end{proof}

\noindent The residual gap $\epsilon > 0$ reflects the fundamental
limits of structural verification under zero-trust: a perfect gap of
zero would require either a trusted verifier or an unbounded
attestation budget. The bound~\eqref{eq:integrity_gap_bound} makes
this tradeoff explicit, with $\kappa$ estimated online, $\rho^*$ set
adaptively, and $b$ depends on the size of agent group, leaving no free parameters
requiring human specification.


\subsection{How the Calibrated Equilibrium Addresses Each Debate Failure (Figure \ref{fig:intro})}
\label{app:cal-equi-adress-mad-failure}

The calibration mechanism $\Phi$ (Section~\ref{sec:calibration}) replaces the vote-based commitment rule $C$ with a score-based rule $\tilde{C}$ (Eq.~\ref{eq:cali-commit-rule} in Appendix~\ref{app:gap_closing}), which removes the dependency of the commit outcome on agent votes and therefore on agent confidence. We trace how this single change resolves each of the three cases.

\emph{Case 1 (confident corruption).} In the uncalibrated game, an over-confident contributor's $\Delta_t$ commits whenever auditors rubber-stamp (Prop.~\ref{prop:freeriding}) or collude (Prop.~\ref{prop:collusion}), regardless of $\Delta_t$'s alignment with $\mathcal{M}_t$. Under $\Phi$, commitment requires $\rho(\Delta_t, \mathcal{M}_t) \geq \rho^*$. A corrupt $\Delta_t$ (whether hallucinated, sycophantic, or adversarially crafted) produces low $\rho_{\mathrm{align}}$ because the entry disrupts retrieval over auditor probes against $\mathcal{M}_t$ (\S\ref{sec:calibration}), and often low $\rho_{\mathrm{detect}}$ because the entry lacks the structural support of genuinely informative content. The contributor's stated confidence does not enter the score. Channel~1 of Appendix~\ref{app:gap_closing} (\emph{Raising the contributor's detection cost}) formalizes this: the manipulation ceiling $\tilde{D}(\Delta_t^*, \mathcal{M}_t) \leq 1/(1+\kappa)$ bounds the detectability of any corrupt $\Delta_t$ that passes $\Phi$, and this bound tightens as $\kappa$ grows.

\emph{Case 2 (confident over-rejection).} In the uncalibrated game, a blocking coalition (Prop.~\ref{prop:adversarial_rejection}) or a confident colluding majority (Prop.~\ref{prop:collusion}) can suppress a tentative but correct $\Delta_t$ simply by voting against it. Under $\Phi$, the auditor's vote does not enter the commit decision at all, only the structural evidence the auditor submits (probe queries for embedding memory, reachability checks for graph memory) contributes to $\rho$. A correct $\Delta_t$ produces non-disrupting retrieval against well-formed probes (high $\rho_{\mathrm{align}}$) and adequate structural support (high $\rho_{\mathrm{detect}}$) regardless of any auditor's stated rejection. Channel~2 of Appendix~\ref{app:gap_closing} (\emph{Eliminating auditor free-riding}) formalizes this: $\partial \mathbb{E}[\mathrm{Align}]/\partial v_i = 0$, so any vote-based suppression strategy has zero leverage on the commit outcome.

\emph{Case 3 (collective hedging).} Cases 1 and 2 are active failures: the wrong outcome is committed or the right outcome is rejected. Case 3 is a passive failure: uncertain (or no) commit happens when a useful one was possible. $\Phi$ does not attempt to ``force'' a commit in Case 3, when all auditors hedge, $\rho_{\mathrm{align}}$ is low across the board because hedging probes produce minimal structural evidence, and the system correctly declines to commit. This is the right behavior: committing a noise-level $\Delta_t$ under collective hedging would be a Case 1 failure in disguise. What $\Phi$ provides in Case 3 is \emph{symmetry of failure}: an honestly-hedged round commits nothing, and an adversarially-hedged round also commits nothing, so an attacker has no asymmetric leverage. Channel~3 of Appendix~\ref{app:gap_closing} (\emph{Deterring collusion intertemporally}) ensures that auditors who systematically submit hedge-style probes accumulate credibility decay, so an adversary cannot maintain Case 3 outcomes indefinitely against an honest minority.

Two observations close this analysis. First, $\Phi$ does not require any agent to be honest: the score $\rho$ is computed entirely from structural properties of $\Delta_t$ and $\mathcal{M}_t$, so the calibrated equilibrium is robust to all three failure modes simultaneously, even when they coexist in the same round (Remark~\ref{rem:failure_modes}). Second, the three cases are addressed by \emph{different channels} of the calibration mechanism (Case 1 by $\rho_{\mathrm{detect}}$ bounding manipulation; Case 2 by $\rho$ being independent of votes; Case 3 by $\rho_{\mathrm{align}}$ declining to commit under hedge), which is why no single component of $\Phi$ is sufficient on its own. The ablation study in Appendix~\ref{app:ablation} confirms that removing any one component degrades performance, and the degradation pattern matches which case the removed component primarily addresses.